\subsection{Project - AI for Medical Diagnosis and Prediction | Week #4: Tumor Classification and Segmentation}

\subsubsection*{Your task}
This week, you will continue the analysis pipeline from Week #3 and add a **segmentation module** to your framework. You will then revisit classification, this time with Convolutional Neural Networks (CNN). The goal is to compare the performance of CNNs with previously trained machine learning models. 

In addition, you will train a segmentation model with a classification head. The objective is to see if the segmentation of relevant areas helps for pathology detection.

In addition, this week will serve as your **MidTerm Evaluation**. You will deliver a report combining insights and results from:
\begin{itemize}
  \item Week 2: Feature extraction and data exploration.
  \item Week 3: Pathology classification using radiomics and machine learning.
  \item Week 4: Pathology classification and segmentation using CNNs.
\end{itemize}

\subsubsection*{Instructions}
\begin{enumerate}
  \item \textbf{Classification Module.}
    \begin{itemize}
      \item Load the dataset using the same train, validation and test splits as during the previous week.
      \item Implement a PyTorch Dataset class for classification.
      \item Train a 2D CNN (e.g., ResNet18, EfficientNet-B0).
      \item Evaluate performance using:
        \begin{itemize}
          \item Accuracy score
          \item F1-score
          \item ROC curve and AUC 
          \item Precision and Recall
        \end{itemize}
      \item Save the model.
    \end{itemize}

  \item \textbf{Segmentation-Guided Classification.}
    \begin{itemize}
      \item Implement a PyTorch Dataset class for segmentation.
      \item Train a 2D segmentation model (e.g., UNet, FCN) or re-use a pre-trained Foundation Model (e.g., MedSAM).
      \item Add a classification head on top of the segmentation model. 
      \item Evaluate segmentation performance using:
        \begin{itemize}
          \item Dice Similarity Coefficient (DSC)
          \item Intersection-over-Union (IoU)
          \item Precision 
        \end{itemize}
      \item Evaluate classification performance using the same metrics as above. 
      \item Save the best model.
      \item Compare the performance to your previous classifier.
    \end{itemize}

  \item \textbf{MidTerm Evaluation Report.}
    \begin{itemize}
      \item Write a concise (max 4 pages) report including:
        \begin{itemize}
          \item Description of dataset and preprocessing pipeline.
          \item Comparison of machine learning models, CNN and classification results.
          \item Plots and tables of metrics and sample outputs.
          \item Reflections: What worked? What didn’t? What would you try next?
        \end{itemize}
      \item Deliver the report as a PDF along with your notebook.
    \end{itemize}
\end{enumerate}

\subsubsection*{Evaluation Questions}
\begin{enumerate}
  \item How can segmentation outputs improve classification tasks?
  \item What are the trade-offs between pixel-level and image-level supervision?
  \item What metric would you use to compare two segmentation models, and why?
  \item If your segmentation model produces false positives, how might that affect downstream classification?
\end{enumerate}

\subsubsection*{Guidelines}
[STANDARD GUIDELINES TO BE ADDED BY LD IN ULTRA - DO NOT AMEND]
